\documentclass[12pt]{article}
\usepackage[landscape,margin=0.4in,bottom=0.55in,top=0.4in]{geometry}
\usepackage[T1]{fontenc}
\usepackage[utf8]{inputenc}
\usepackage[spanish]{babel}
\usepackage{enumitem}
\usepackage{xcolor}
\definecolor{solutionred}{HTML}{7F1D1D}
\newcommand{\sol}[1]{{\color{solutionred}\small\strut #1\strut}}
\setlength{\parindent}{0pt}
\setlist[enumerate]{leftmargin=0.65cm,topsep=3mm,parsep=0pt,itemsep=3mm}
\usepackage{array}
\usepackage{longtable}
\newcolumntype{L}[1]{>{\raggedright\arraybackslash}p{#1}}
\begin{document}
\raggedbottom
\color{black}
\begin{center}
{\Large\bfseries C3--C4: ESCUCHAR A LOS USUARIOS}\\[1mm]
Emprendimiento e Innovación Grado 10\\[1mm]
\textbf{Ejemplo de encuestas completas}
\end{center}
\textbf{Simulación pedagógica:} todos los estudiantes y sus respuestas son ficticios. Este documento no contiene datos de encuestas realizadas en el Colegio San Bartolomé de la Merced.

\smallskip
El caso corresponde a la espera para comprar alimentos en el punto de venta de la cafetería. Los usuarios A, B y C conservan los perfiles del documento de solución; D y E amplían la variedad de experiencias simuladas. Los tiempos expresados son estimaciones de los personajes, no mediciones verificadas. Las respuestas ejemplifican el registro de experiencias para C3 y C4; no representan al conjunto del colegio.

\section{Cuestionario de la encuesta}
\begin{enumerate}
\item En tu último descanso, ¿qué hiciste para conseguir alimentos y en qué momento?
\item Si esperaste para comprar, ¿cuánto calculas que esperaste antes de la atención?
\item ¿Cómo decidiste qué comprar y qué información consultaste?
\item ¿Qué ocurrió desde que llegaste a la fila hasta que recibiste tu compra o te retiraste?
\item Después de comprar o retirarte, ¿qué pudiste hacer en el tiempo restante del descanso?
\item ¿Cómo supiste dónde hacer la fila y cuándo era tu turno?
\item En esa experiencia, ¿qué funcionó bien y qué dificultades, si hubo alguna, encontraste?
\end{enumerate}
\newpage
\section{Respuestas de la encuesta}
Ejemplos ficticios de cinco estudiantes anónimos. Cada fila corresponde a una pregunta del cuestionario; las respuestas conservan su redacción original.
\begingroup\small
% Evita contracción infinita al dividir la tabla entre páginas.
\def\vss{\vskip 0pt plus 1fil\relax}
\setlength{\tabcolsep}{4pt}\renewcommand{\arraystretch}{1.12}
\setlength{\LTleft}{0pt plus 1fill}\setlength{\LTright}{0pt plus 1fill}
\begin{longtable}{|*{6}{L{\dimexpr(\textwidth-12\tabcolsep-7\arrayrulewidth)/6\relax}|}}
\hline \textbf{Pregunta} & \textbf{Usuario A} & \textbf{Usuario B} & \textbf{Usuario C} & \textbf{Usuario D} & \textbf{Usuario E} \\ \hline
\endfirsthead
\hline \textbf{Pregunta} & \textbf{Usuario A} & \textbf{Usuario B} & \textbf{Usuario C} & \textbf{Usuario D} & \textbf{Usuario E} \\ \hline
\endhead\endfoot\endlastfoot
\textbf{1.} En tu último descanso, ¿qué hiciste para conseguir alimentos y en qué momento? & \sol{Fui al punto de venta apenas empezó el descanso porque no había traído comida. Quería comprar algo para comer antes de volver a clase.} & \sol{Primero conversé con unos compañeros. Fui a comprar cuando ya había pasado una parte del descanso y vi menos gente en el punto de venta.} & \sol{Pasé primero por el baño y después fui al punto de venta para comprar algo de comer. Ya había transcurrido parte del descanso cuando me uní a la fila.} & \sol{Al comenzar el descanso comí el sándwich que había traído de casa. Más tarde fui al punto de venta porque quería comprar agua.} & \sol{Fui al punto de venta poco después de empezar el descanso. Llevaba dinero para una sola cosa y quería comprar algo que me quitara el hambre.} \\ \hline
\textbf{2.} Si esperaste para comprar, ¿cuánto calculas que esperaste antes de la atención? & \sol{Calculo que esperé unos siete minutos antes de que me atendieran; no lo cronometré.} & \sol{Me pareció que fueron unos tres minutos antes de la atención. Es un cálculo de memoria, no una medición.} & \sol{Esperé un rato, pero no sabría decir cuántos minutos. Me retiré antes de recibir atención, así que no puedo estimar una espera completa hasta ser atendido.} & \sol{Calculo que esperé alrededor de dos minutos antes de que me atendieran. No miré el reloj para medirlo.} & \sol{Creo que esperé unos cinco minutos antes de la atención. Es una estimación; no registré la hora de llegada ni la de mi turno.} \\ \hline
\textbf{3.} ¿Cómo decidiste qué comprar y qué información consultaste? & \sol{Pensaba comprar una empanada y una bebida. No sabía el precio de la empanada y lo pregunté cuando me atendieron; al conocer el total, decidí comprar las dos cosas.} & \sol{Ya sabía que quería un pan y una bebida que había comprado antes. Recordaba sus precios; al pedir confirmé que el total seguía siendo el que esperaba.} & \sol{Tenía en mente comprar un pan porque quería algo fácil de comer. No alcancé a consultar el precio ni a confirmar si había; me fui antes de hacer el pedido.} & \sol{Solo quería agua, así que no comparé alimentos. Miré las bebidas que alcanzaba a ver y pregunté el precio de la botella al hacer el pedido para saber si me alcanzaba el dinero.} & \sol{Pensaba comprar una empanada. Al llegar pregunté por ella y me dijeron que en ese momento no había. Pregunté qué panes tenían y cuánto costaban; elegí uno que podía pagar con lo que llevaba.} \\ \hline

\textbf{4.} ¿Qué ocurrió desde que llegaste a la fila hasta que recibiste tu compra o te retiraste? & \sol{Me ubiqué detrás del último estudiante y avancé con la fila. Al llegar mi turno pregunté el precio, hice el pedido y pagué. Me entregaron la empanada y la bebida.} & \sol{Llegué al final de la fila, esperé mientras avanzaba y pedí lo que tenía pensado cuando me tocó. Pagué y recibí ambas cosas sin tener que cambiar mi pedido.} & \sol{Me puse al final y avancé un poco. Vi que todavía tenía estudiantes delante y no sabía cuánto faltaba para llegar. Como sentía que me quedaba poco descanso, salí de la fila sin comprar.} & \sol{Me uní a una fila corta. Cuando me atendieron pedí agua y confirmé el precio. La persona fue a buscar más botellas para entregarme una; pagué y la recibí. No sé cuánto duró ese momento de reposición.} & \sol{Esperé detrás de otros estudiantes. Cuando me tocó, consulté por la empanada, escuché las opciones disponibles y cambié mi elección por un pan. Pagué y me lo entregaron.} \\ \hline
\textbf{5.} Después de comprar o retirarte, ¿qué pudiste hacer en el tiempo restante del descanso? & \sol{Comí de afán y tomé la bebida. Casi no alcancé a conversar con mis amigos antes de regresar a clase.} & \sol{Me senté cerca con mis compañeros y pude comer tranquilo. Después regresamos juntos al salón.} & \sol{Fui a tomar agua y regresé hacia el salón. No comí en ese descanso y seguí con hambre hasta que tuve otra oportunidad de comer.} & \sol{Tomé agua y me quedé conversando un rato. Como ya había comido antes de comprar, no tuve que apurarme para terminar la comida.} & \sol{Alcancé a comer el pan y a conversar brevemente con una amiga. Me hubiera gustado seguir hablando, pero ya teníamos que volver a clase.} \\ \hline
\textbf{6.} ¿Cómo supiste dónde hacer la fila y cuándo era tu turno? & \sol{Vi dónde terminaba la fila y me puse allí. Reconocí a quien estaba delante de mí; cuando terminó su compra y la persona que atendía me llamó, supe que seguía yo.} & \sol{La fila se veía desde donde llegué. Me puse detrás del último estudiante y esperé a que atendieran a quienes estaban delante; luego me indicaron que podía pedir.} & \sol{Encontré el final mirando a los estudiantes que esperaban y pregunté al último si estaba en la fila. Sabía detrás de quién iba, pero me retiré antes de que llegara mi turno.} & \sol{Vi a dos estudiantes esperando uno detrás del otro y me ubiqué al final. Cuando salió quien estaba delante, la persona que atendía me preguntó qué necesitaba.} & \sol{Me ubiqué detrás de la última persona. Algunas personas pasaron al lado de la fila, pero seguí pendiente del estudiante que iba delante y no perdí mi turno.} \\ \hline
\textbf{7.} En esa experiencia, ¿qué funcionó bien y qué dificultades, si hubo alguna, encontraste? & \sol{Recibí lo que pedí y tuve claro mi turno. La dificultad fue que solo confirmé el precio al llegar adelante y después me quedó poco tiempo para comer con calma.} & \sol{La fila avanzó y pude comprar lo que quería. En esa ocasión no encontré una dificultad; como ya tenía decidido el pedido, la compra me resultó sencilla.} & \sol{Pude ubicarme en la fila sin problema. La dificultad fue no poder calcular si alcanzaría a comprar y comer; no tuve una experiencia de pago o entrega porque abandoné la compra.} & \sol{Conseguí la bebida y pude pagar con el dinero que llevaba. Tuve que esperar un poco mientras buscaban más agua, pero en esta ocasión eso no me impidió aprovechar el descanso.} & \sol{Me explicaron las opciones y encontré algo dentro de mi presupuesto. La dificultad fue enterarme al pedir de que no había lo que quería; tuve que decidir otra cosa en ese momento.} \\ \hline
\end{longtable}\endgroup
\end{document}






